\subsection{Awareness and Cognitive Biases: A Unifying Framework}
\label{subsec:11-9-biases}

%%%%%%%%%%%%%%%%%%%%%%%%%%%%%%%%%%%%%%%%%%%%%%%%%%%%%%%%%%%%%%%%%%%%%%%%%%%%%%%
% SECTION 11.9: Awareness and Cognitive Biases
% Narrative overview - Formal specifications in Appendix AU Corollaries
%%%%%%%%%%%%%%%%%%%%%%%%%%%%%%%%%%%%%%%%%%%%%%%%%%%%%%%%%%%%%%%%%%%%%%%%%%%%%%%

A remarkable property of the Awareness function is its capacity to explain
and systematically organize a wide range of cognitive biases. Rather than
treating each bias as an isolated phenomenon, the Awareness framework reveals
biases as \emph{derived consequences} of the axiomatic structure.

%------------------------------------------------------------------------------
\subsubsection{The Core Insight: Biases as Corollaries}
%------------------------------------------------------------------------------

\begin{tcolorbox}[colback=blue!5!white, colframe=blue!75!black, title=Central Thesis]
\textbf{Cognitive biases are not random errors or independent phenomena. They are formal corollaries---theorems derived from axiom combinations:}
\begin{equation}
\text{Bias}_i = f(\text{AWX-A}_{j_1}, \text{AWX-A}_{j_2}, \ldots, \text{AWX-A}_{j_k})
\end{equation}
Each bias has a specific derivation from one or more AWX axioms.
\end{tcolorbox}

This corollary structure has three implications:
\begin{enumerate}
\item \textbf{Explanatory}: Each bias traces to specific axiom combinations
\item \textbf{Predictive}: We can predict when biases occur based on which axioms are active
\item \textbf{Interventional}: To reduce bias $B$, intervene on its constituent axioms
\end{enumerate}

%------------------------------------------------------------------------------
\subsubsection{Eight Categories of Derived Biases}
%------------------------------------------------------------------------------

The 59 formally derived corollaries (COR-01 to COR-59) are organized into eight categories:

\begin{table}[htbp]
\centering
\caption{Cognitive Biases as Derived Corollaries}
\label{tab:bias-corollaries-summary}
\renewcommand{\arraystretch}{1.3}
\begin{tabular}{@{}llll@{}}
\toprule
\textbf{Category} & \textbf{Corollaries} & \textbf{Primary Axioms} & \textbf{Example Bias} \\
\midrule
System-Level & COR-01--07 & A4, A35, A58 & Present Bias, Akrasia \\
Self-Community & COR-08--15 & A1a, A1b, A34 & Egocentric Bias, Bystander Effect \\
Utility-Type & COR-16--23 & A34--A37 & Tragedy of Commons, Free Rider \\
Belief-Based & COR-24--33 & A81--A90 & Confirmation Bias, Overconfidence \\
Cognitive & COR-34--41 & A9, A14, A27 & Availability, Anchoring \\
Emotional & COR-42--47 & A15--A17 & Loss Aversion, Negativity Bias \\
Social & COR-48--53 & A19--A21, A70 & Conformity, Groupthink \\
Temporal & COR-54--59 & A58--A60, A63 & Recency, Status Quo Bias \\
\midrule
\textbf{Total} & \textbf{59} & & \\
\bottomrule
\end{tabular}
\end{table}

%------------------------------------------------------------------------------
\subsubsection{Key Theoretical Insights}
%------------------------------------------------------------------------------

\paragraph{Axiom A37 is Pivotal.}
The single axiom ``KNU is never instant'' (AWX-A37) generates six corollaries:
\begin{itemize}
\item COR-05 (Temptation): $A_{INU}^{instant} = 1 \land A_{KNU}^{instant} = 0$
\item COR-14 (Empathy Gap): $A^{comm}$ requires deliberate effort
\item COR-16 (Tragedy of Commons): Individual benefits instant; collective costs not
\item COR-17 (Free Rider): $A_{INU} > A_{KNU}$ structurally
\item COR-18 (Social Dilemmas): INU crowds out KNU
\item COR-19 (NIMBY): Self-INU instant; Community-KNU not
\end{itemize}
This explains why collective action problems are so persistent across cultures and contexts.

\paragraph{Belief System Generates Motivated Cognition.}
Ten corollaries (17\% of total) derive primarily from the Belief System (A81--A90):
\begin{equation}
\text{COR-24 to COR-33} \Leftarrow \{A81, A82, A84, A85, A89\}
\end{equation}
This includes confirmation bias, denial, rationalization, overconfidence, and optimism bias.

\paragraph{Loss Aversion Has Micro-Foundations.}
COR-43 (Loss Aversion) derives from A16 (Fear) and A8 (Asymmetry):
\begin{equation}
|A_{fear}| > |A_{hope}| \Rightarrow U(x) - U(0) < |U(-x) - U(0)| \Rightarrow \lambda > 1
\end{equation}
This provides a micro-foundation for Prospect Theory's loss aversion parameter.

%------------------------------------------------------------------------------
\subsubsection{Intervention Logic}
%------------------------------------------------------------------------------

The corollary structure enables systematic debiasing:

\begin{table}[htbp]
\centering
\caption{Debiasing Through Axiom Intervention}
\label{tab:debiasing-axiom}
\renewcommand{\arraystretch}{1.3}
\begin{tabular}{@{}lll@{}}
\toprule
\textbf{Bias Category} & \textbf{Target Axioms} & \textbf{Intervention Strategy} \\
\midrule
System-Level & A4, A35 & Cooling-off periods; pre-commitment \\
Self-Community & A1b, A37 & Identifiable victims; perspective-taking \\
Utility-Type & A37, A42 & Make KNU visible; align IDN-KNU \\
Belief-Based & A82, A85, A89 & Increase accuracy stakes; external accountability \\
Cognitive & A14, A9 & Statistical formats; slow processing \\
Emotional & A16, A17 & Balanced emotional framing \\
Social & A19, A70 & Diverse exposure; dissent protection \\
Temporal & A58, A60 & Varied triggers; future salience \\
\bottomrule
\end{tabular}
\end{table}

%------------------------------------------------------------------------------
\subsubsection{Extensibility}
%------------------------------------------------------------------------------

The corollary framework is \textbf{open}: additional biases can be derived as new combinations of existing axioms. For example:

\begin{itemize}
\item \textbf{Dunning-Kruger}: COR-30 $\land$ A51 $\land$ A74
\item \textbf{Planning Fallacy}: COR-31 $\land$ COR-06 $\land$ A46
\item \textbf{Sunk Cost}: COR-28 $\land$ A85 $\land$ COR-43
\end{itemize}

\begin{cross-reference}
\textbf{Formal Specification:} All 59 corollaries are formally derived in Appendix~\ref{app:bcm-axiom-formalization}, Section ``Corollaries: Cognitive Biases as Derived Phenomena.'' Each corollary specifies: ID, Bias Name, Derivation (axiom combination), Formal Statement, and Intervention Logic.
\end{cross-reference}

%%%%%%%%%%%%%%%%%%%%%%%%%%%%%%%%%%%%%%%%%%%%%%%%%%%%%%%%%%%%%%%%%%%%%%%%%%%%%%%
% END OF SECTION 11.9
%%%%%%%%%%%%%%%%%%%%%%%%%%%%%%%%%%%%%%%%%%%%%%%%%%%%%%%%%%%%%%%%%%%%%%%%%%%%%%%
